% !TeX root = ../../main.tex
% =====================================================================
%  El neumático: origen de todas las fuerzas
%  Responsable:
%  Última revisión:
% =====================================================================
\chapter{El neumático: origen de todas las fuerzas}
\label{ch:neumatico}

% Las curvas de este capítulo usan la función mf(x,B,C,D,E) (Magic Formula),
% declarada una sola vez en Estilo/preambulo.tex.

\begin{objetivos}
  \item Explicar de dónde sale el agarre de un neumático y por qué su
        coeficiente puede pasar de 1, cosa que el rozamiento de Coulomb prohíbe.
  \item Definir ángulo de deriva y deslizamiento longitudinal con el convenio de
        signos del equipo, y leer una curva $\Fy(\slipa)$: zona lineal, pico y
        caída.
  \item Usar la elipse de fricción para calcular cuánta fuerza lateral queda
        disponible mientras se frena.
  \item Cuantificar la sensibilidad a la carga vertical y estimar cuánto agarre
        pierde un eje por transferencia de carga.
  \item Explicar el efecto de caída, presión y temperatura, y saber diagnosticar
        los tres con un pirómetro.
\end{objetivos}

\fichasesion{90 minutos}{Capítulos \ref{ch:coche_como_sistema} y
\ref{ch:mecanica_aplicada}}%
{Manómetro de precisión, pirómetro de aguja, un neumático desmontado y, si se
tienen, las curvas del neumático propio}

Todo lo que hace el coche ---acelerar, frenar, girar--- lo hacen cuatro huellas
de contacto del tamaño de la palma de una mano. El motor no mueve el coche: el
neumático lo hace, y el motor solo le da un par. La aerodinámica no gira el
coche: cambia la carga vertical del neumático para que él gire más. Por eso
este capítulo abre el bloque de dinámica vehicular, y por eso todas las
decisiones de suspensión, aerodinámica y transmisión de los bloques siguientes
se justifican, en última instancia, por lo que le hacen al neumático.

\begin{clave}
\textbf{Las curvas de este capítulo son sintéticas.} Están generadas con un juego
de coeficientes de Pacejka plausible para un neumático de Formula Student, no son
datos medidos de ningún neumático concreto. Sirven para explicar \emph{la forma}
de las curvas y los órdenes de magnitud, nunca para diseñar. Los números con los
que se diseña salen del capítulo \ref{ch:modelado_neumatico}.
\end{clave}

\pendiente{Sustituir las curvas sintéticas por las del neumático que montamos en
cuanto existan, aunque sean de un neumático parecido del TTC. Mientras tanto, la
advertencia de arriba se queda.}

% ---------------------------------------------------------------------
\section{Mecanismos de agarre}
\label{sec:neumatico:01}
\index{agarre}\index{adhesión}\index{histéresis}

\subsection{Por qué el rozamiento de Coulomb no sirve aquí}

En primero se aprende que $F = \mu N$, que $\mu$ es una constante del par de
materiales, que no depende del área de contacto y que vale menos que 1. Las tres
cosas son falsas para un neumático:

\begin{itemize}
  \item $\mu$ \textbf{no es constante}: baja al aumentar la carga vertical
        (apartado \ref{sec:neumatico:04}), depende de la temperatura, de la
        presión, de la velocidad de deslizamiento y del acabado del asfalto.
  \item $\mu$ \textbf{sí depende del área}: por eso un neumático ancho agarra
        más, aunque la carga sea la misma.
  \item $\mu$ \textbf{puede pasar de 1} con holgura. Un neumático de competición
        blando da entre \num{1.4} y \num{1.8} en pista, y sobre la banda de
        ensayo de un laboratorio llega a \numrange{2.4}{2.9}.
\end{itemize}

La razón es que la goma no roza: se deforma. El agarre de un neumático nace de
dos mecanismos distintos que trabajan a la vez.

\subsection{Adhesión e histéresis}

\begin{description}[leftmargin=0pt,labelsep=0pt,style=nextline]
  \item[Adhesión] Enlaces moleculares que se forman y se rompen continuamente
    entre la goma y la superficie. Es el mecanismo dominante sobre asfalto
    liso y limpio, y es el que \emph{escala con el área real de contacto}: de
    ahí que la anchura del neumático importe. También es el que desaparece
    cuando hay agua, polvo o goma vieja entre medias.
  \item[Histéresis] La goma se mete en los huecos entre los granos del asfalto,
    se comprime por delante de cada grano y no recupera del todo por detrás.
    Esa asimetría es una fuerza neta que se opone al deslizamiento. Es el
    mecanismo dominante sobre superficies rugosas y el que sigue funcionando
    en mojado, porque no necesita contacto molecular limpio.
\end{description}

Hay dos más ---desgaste (arranque de material) y cohesión--- pero en un coche
que funciona bien aportan poco: si el neumático agarra por desgaste, se está
destrozando.

\begin{clave}
La goma es un material \textbf{viscoelástico}: su rigidez y su amortiguamiento
dependen de la temperatura y de la frecuencia a la que se la excita. De ahí salen
dos consecuencias que gobiernan todo lo demás:

\begin{enumerate}
  \item El agarre tiene una \textbf{ventana de temperatura}. Frío o caliente de
        más, el mismo neumático da bastante menos.
  \item El agarre depende de la \textbf{velocidad de deslizamiento}, es decir, de
        cuánto y cómo se está deslizando la huella. Por eso el neumático necesita
        deslizar un poco para dar su máximo, que es justo lo que explica el
        apartado siguiente.
\end{enumerate}
\end{clave}

\subsection{Órdenes de magnitud}

La huella de contacto se estima, en primera aproximación, dividiendo la carga
vertical por la presión de inflado:

\begin{equation}
  A_{\text{huella}} \approx \frac{\Fz}{\pneu}
  \label{eq:area_huella}
\end{equation}

Con \SI{800}{\newton} sobre la rueda y \SI{0.83}{\bar} de presión salen unos
\SI{96}{\centi\meter\squared}, es decir, algo así como \num{12} por
\SI{8}{\centi\meter}. La estimación se queda corta porque parte de la carga la
lleva la carcasa y no el aire, pero da el orden de magnitud correcto y sirve para
discutir: \textbf{el coche entero apoya en cuatro rectángulos que caben en un
folio}.

\begin{center}
\begin{tabularx}{\textwidth}{@{}lXl@{}}
\toprule
\textbf{Superficie} & \textbf{Situación} & \textbf{$\muy$ típico} \\
\midrule
Asfalto seco, neumático de calle & Coche de serie & \numrange{0.8}{1.0} \\
Asfalto seco, neumático semi-slick & Neumático duro de FS & \numrange{1.1}{1.4} \\
Asfalto seco, slick de competición & Neumático blando de FS & \numrange{1.4}{1.8} \\
Banda de ensayo (3M Safety Walk) & Datos brutos del TTC & \numrange{2.4}{2.9} \\
Asfalto mojado & Cualquiera & la mitad, o menos \\
\bottomrule
\end{tabularx}
\end{center}

La última fila de la tabla es la que más disgustos da: los datos de laboratorio
no son datos de pista, y usarlos sin escalar produce un coche diseñado para un
agarre que nunca va a tener. El capítulo \ref{ch:modelado_neumatico} se ocupa de
eso en detalle.

% ---------------------------------------------------------------------
\section{Ángulo de deriva y deslizamiento longitudinal}
\label{sec:neumatico:02}
\index{ángulo de deriva}\index{deslizamiento longitudinal}\index{modelo de cepillo}

\subsection{Definiciones y convenio de signos}

El \textbf{ángulo de deriva} $\slipa$ es el ángulo entre la dirección en la que
apunta la rueda y la dirección en la que realmente se mueve:

\begin{equation}
  \slipa = \arctan\!\left(\frac{v_y}{v_x}\right)
  \label{eq:slip_angle}
\end{equation}

El \textbf{deslizamiento longitudinal} $\slipr$ compara la velocidad periférica
de la rueda con la velocidad de avance:

\begin{equation}
  \slipr = \frac{\omega\, r_e - v_x}{|v_x|}
  \label{eq:slip_ratio}
\end{equation}

donde $r_e$ es el radio efectivo de rodadura. Con este convenio, $\slipr > 0$ es
tracción (la rueda gira más deprisa que el suelo), $\slipr < 0$ es frenada, y
$\slipr = -1$ es la rueda bloqueada.

\begin{ojo}
El convenio de signos es la primera fuente de errores del bloque, y siempre
aparece tarde: en la simulación de vuelta, cuando el coche gira al revés.
Fíjalo una vez ---sistema SAE, según la nomenclatura de este libro--- y
compruébalo en cada modelo que uses: los datos del TTC, la herramienta de
cinemática y el simulador no tienen por qué coincidir entre sí. El error clásico
es un $\Fy$ con el signo cambiado respecto de $\slipa$.
\end{ojo}

\subsection{Qué pasa dentro de la huella}

La imagen mental correcta es el \textbf{modelo de cepillo}: la banda de rodadura
es una fila de cerdas elásticas que entran en la huella por delante y salen por
detrás. Una cerda entra sin deformar y, mientras avanza por la huella, el suelo
la arrastra lateralmente; se va estirando y genera fuerza. Puede seguir
estirándose mientras la fuerza elástica no supere el límite de rozamiento local,
que es proporcional a la presión de contacto: alta en el centro de la huella y
nula en los bordes. Cuando lo supera, la cerda desliza.

De ahí sale toda la forma de la curva (\cref{fig:neumatico:cepillo}):

\begin{itemize}
  \item Con $\slipa$ pequeño casi toda la huella está \textbf{adherida}: la
        fuerza crece proporcionalmente al ángulo. Es la zona lineal.
  \item Al aumentar $\slipa$, la zona de \textbf{deslizamiento} crece desde la
        parte trasera de la huella hacia delante.
  \item El \textbf{pico} de fuerza se alcanza cuando una fracción grande de la
        huella ya desliza, no cuando toda está adherida.
  \item Pasado el pico, casi toda la huella desliza, la fuerza cae y el
        neumático se calienta y se desgasta muy deprisa.
\end{itemize}

\begin{figure}[htbp]
\centering
\begin{tikzpicture}[x=1cm,y=1cm]

% ---- Panel izquierdo: ángulo de deriva pequeño ----
\begin{scope}
  \fill[fusalPrim!6] (0,-1.1) rectangle (6,1.1);
  \draw[fusalGris] (0,-1.1) rectangle (6,1.1);
  % límite de rozamiento (parábola de presión)
  \draw[fusalGris,densely dashed]
       plot[domain=0:6,samples=40] (\x,{0.9*\x*(6-\x)/9});
  % zona adherida
  \fill[fusalPrim!25] (0,0) -- plot[domain=0:4.8,samples=2] (\x,{0.12*\x})
       -- (4.8,0) -- cycle;
  % zona deslizante
  \fill[fusalSec!25] (4.8,0) -- plot[domain=4.8:6,samples=30]
       (\x,{0.9*\x*(6-\x)/9}) -- (6,0) -- cycle;
  \draw[fusalPrim,very thick] plot[domain=0:4.8,samples=2] (\x,{0.12*\x});
  \draw[fusalSec,very thick] plot[domain=4.8:6,samples=30]
       (\x,{0.9*\x*(6-\x)/9});
  \draw[fusalGris,-{Latex[length=1.6mm]}] (0,0) -- (6.4,0);
  \node[cursonota,anchor=north] at (3,-1.25) {huella, $\slipa$ pequeño};
  \node[cursonota,anchor=south,text=fusalPrim] at (2.2,0.32) {adherida};
  \node[cursonota,anchor=north east,text=fusalSec] at (5.9,-0.15) {desliza};
  \node[cursonota,anchor=east] at (-0.15,0) {entra};
  \node[cursonota,anchor=west] at (6.55,0) {sale};
\end{scope}

% ---- Panel derecho: ángulo de deriva grande ----
\begin{scope}[xshift=8.4cm]
  \fill[fusalPrim!6] (0,-1.1) rectangle (6,1.1);
  \draw[fusalGris] (0,-1.1) rectangle (6,1.1);
  \draw[fusalGris,densely dashed]
       plot[domain=0:6,samples=40] (\x,{0.9*\x*(6-\x)/9});
  \fill[fusalPrim!25] (0,0) -- plot[domain=0:1.0,samples=2] (\x,{0.5*\x})
       -- (1.0,0) -- cycle;
  \fill[fusalSec!25] (1.0,0) -- plot[domain=1.0:6,samples=40]
       (\x,{0.9*\x*(6-\x)/9}) -- (6,0) -- cycle;
  \draw[fusalPrim,very thick] plot[domain=0:1.0,samples=2] (\x,{0.5*\x});
  \draw[fusalSec,very thick] plot[domain=1.0:6,samples=40]
       (\x,{0.9*\x*(6-\x)/9});
  \draw[fusalGris,-{Latex[length=1.6mm]}] (0,0) -- (6.4,0);
  \node[cursonota,anchor=north] at (3,-1.25) {huella, $\slipa$ grande};
  \node[cursonota,anchor=south,text=fusalPrim] at (0.9,0.5) {adh.};
  \node[cursonota,anchor=south,text=fusalSec] at (3.6,0.55) {desliza};
\end{scope}

\end{tikzpicture}
\caption[Modelo de cepillo: zonas de adherencia y deslizamiento]{Modelo de
cepillo. La línea gruesa es la deformación lateral de la banda de rodadura a lo
largo de la huella; la discontinua gris, el límite que impone el rozamiento
local (proporcional a la presión de contacto, nula en los bordes). Al crecer el
ángulo de deriva la zona de deslizamiento avanza desde la salida hacia la
entrada. El pico de fuerza lateral se produce con buena parte de la huella ya
deslizando.}
\label{fig:neumatico:cepillo}
\end{figure}

\begin{clave}
\textbf{No hay fuerza sin deslizamiento.} Un neumático que rueda perfectamente
alineado no genera fuerza lateral ninguna. Para que empuje hay que deformarlo, y
deformarlo significa que la rueda apunta a un sitio y se mueve hacia otro. Lo
mismo en longitudinal: para acelerar o frenar, la rueda tiene que girar a una
velocidad distinta de la que le correspondería. Un neumático de Formula Student
da su máximo lateral entre \num{5} y \SI{8}{\degree} de deriva, y su máximo
longitudinal entre el \num{5} y el \SI{15}{\percent} de deslizamiento.
\end{clave}

\subsection{Las dos curvas fundamentales}

La \cref{fig:neumatico:curvas} recoge las dos curvas que hay que tener en la
cabeza. La pendiente de $\Fy(\slipa)$ en el origen es la \textbf{rigidez de
deriva}:

\begin{equation}
  \Calpha = \left.\pderiv{\Fy}{\slipa}\right|_{\slipa=0}
  \label{eq:cornering_stiffness}
\end{equation}

Es el parámetro que gobierna el comportamiento del coche lejos del límite ---o
sea, el \SI{95}{\percent} del tiempo que el piloto pasa conduciendo--- y el que
entra en el gradiente de subviraje del capítulo \ref{ch:regimen_estacionario}.
Un neumático de Formula Student anda entre \num{500} y
\SI{1400}{\newton\per\degree} según la carga.

\begin{figure}[htbp]
\centering
\begin{tikzpicture}
\begin{axis}[cursopanel,
  xlabel={$\slipa$ (\si{\degree})}, ylabel={$\Fy$ (\si{\newton})},
  xmin=0, xmax=13, ymin=0, ymax=3000,
  xtick={0,2,...,12}, ytick={0,500,...,3000},
  title={(a) Fuerza lateral},
  legend pos=south east]
\addplot[curvaA,domain=0:13,samples=140]{mf(x,0.3448,1.5,880,0.3)};
  \addlegendentry{$\Fz=\SI{500}{\newton}$}
\addplot[curvaB,domain=0:13,samples=140]{mf(x,0.3226,1.5,1620,0.3)};
  \addlegendentry{$\Fz=\SI{1000}{\newton}$}
\addplot[curvaC,domain=0:13,samples=140]{mf(x,0.3030,1.5,2220,0.3)};
  \addlegendentry{$\Fz=\SI{1500}{\newton}$}
\addplot[curvaD,domain=0:13,samples=140]{mf(x,0.2857,1.5,2680,0.3)};
  \addlegendentry{$\Fz=\SI{2000}{\newton}$}
% tangente en el origen: rigidez de deriva de la curva de 1000 N
\addplot[fusalGris,densely dotted,domain=0:3.2,samples=2]{784*x};
\node[cursonota,anchor=south east] at (axis cs:3.4,2380) {pendiente $\Calpha$};
\end{axis}
\end{tikzpicture}%
\hfill
\begin{tikzpicture}
\begin{axis}[cursopanel,
  xlabel={$\slipr$ (\si{\percent})}, ylabel={$\Fx$ (\si{\newton})},
  xmin=0, xmax=30, ymin=0, ymax=3000,
  xtick={0,5,...,30}, ytick={0,500,...,3000},
  title={(b) Fuerza longitudinal}]
\addplot[curvaA,domain=0:30,samples=140]{mf(x,0.1567,1.65,930,0.3)};
\addplot[curvaB,domain=0:30,samples=140]{mf(x,0.1425,1.65,1700,0.3)};
\addplot[curvaC,domain=0:30,samples=140]{mf(x,0.1306,1.65,2320,0.3)};
\addplot[curvaD,domain=0:30,samples=140]{mf(x,0.1200,1.65,2800,0.3)};
\node[cursonota,anchor=north east] at (axis cs:29,900)
     {mismas cargas y\\mismos colores que (a)};
\end{axis}
\end{tikzpicture}
\caption[Curvas de fuerza lateral y longitudinal]{Las dos curvas fundamentales
del neumático, para cuatro cargas verticales. (a) Fuerza lateral frente a ángulo
de deriva; la recta punteada es la tangente en el origen, cuya pendiente es la
rigidez de deriva $\Calpha$. (b) Fuerza longitudinal frente a deslizamiento; el
pico aparece en torno al \SI{10}{\percent}, y en frenada la curva es la simétrica
respecto del origen. Curvas sintéticas: sirven para la forma, no como dato de
diseño.}
\label{fig:neumatico:curvas}
\end{figure}

\subsection{Par autoalineante: por dónde se entera el piloto}

La resultante lateral no se aplica en el centro de la huella, sino algo por
detrás, porque la deformación es mayor en la parte trasera. Esa distancia es el
\textbf{avance neumático} $\tneu$, y el momento que produce es el par
autoalineante:

\begin{equation}
  \Mz = -\Fy\,(\tneu + t_m)
  \label{eq:mz}
\end{equation}

donde $t_m$ es el avance mecánico que da la geometría de la dirección
(capítulo \ref{ch:direccion}). Lo interesante es cómo evolucionan las dos
magnitudes (\cref{fig:neumatico:mz}): a medida que la parte trasera de la huella
empieza a deslizar, el centro de presiones se adelanta, $\tneu$ cae y llega a
anularse justo cuando $\Fy$ alcanza su máximo.

\begin{clave}
$\Mz$ hace pico \textbf{mucho antes} que $\Fy$, y se desploma justo cuando el
neumático llega a su máximo. Eso es exactamente lo que el piloto nota como
«se me ha ido el volante ligero»: es el aviso natural de que se está en el
límite. Cualquier decisión de dirección que borre esa señal ---demasiado avance
mecánico, demasiada desmultiplicación, una servoasistencia mal ajustada--- le
quita al piloto el único sensor de límite que lleva puesto.
\end{clave}

\begin{figure}[htbp]
\centering
\begin{tikzpicture}
\begin{axis}[cursopanel,
  xlabel={$\slipa$ (\si{\degree})}, ylabel={fracción del máximo},
  xmin=0, xmax=10, ymin=0, ymax=1.1,
  xtick={0,2,...,10}, ytick={0,0.2,...,1.0},
  title={(a) Fuerza y par, normalizados},
  legend pos=south east]
% modelo de cepillo: lambda = alpha/10 (10 grados = deslizamiento total)
\addplot[curvaA,domain=0:10,samples=120]
   {3*(x/10) - 3*(x/10)^2 + (x/10)^3};
  \addlegendentry{$\Fy/\muy\Fz$}
\addplot[curvaB,domain=0:10,samples=120]
   {(x/10)*(1-x/10)^3/0.10547};
  \addlegendentry{$\Mz/\Mz^{\max}$}
\draw[cursomarca] (axis cs:2.5,0) -- (axis cs:2.5,1.0);
\node[cursonota,anchor=west] at (axis cs:2.7,0.28)
     {pico de $\Mz$};
\end{axis}
\end{tikzpicture}%
\hfill
\begin{tikzpicture}
\begin{axis}[cursopanel,
  xlabel={$\slipa$ (\si{\degree})}, ylabel={$\tneu$ (\si{\milli\meter})},
  xmin=0, xmax=10, ymin=0, ymax=30,
  xtick={0,2,...,10}, ytick={0,5,...,30},
  title={(b) Avance neumático}]
\addplot[curvaC,domain=0:9.98,samples=120]
   {75*(1-x/10)^3/(3 - 3*(x/10) + (x/10)^2)};
\node[cursonota,anchor=north east] at (axis cs:9.7,18)
     {$\tneu \to 0$ en el\\pico de $\Fy$};
\end{axis}
\end{tikzpicture}
\caption[Par autoalineante y avance neumático]{Par autoalineante y avance
neumático según el modelo de cepillo, con deslizamiento total a
\SI{10}{\degree}. (a) El par autoalineante hace pico en torno a un cuarto del
ángulo de deslizamiento total, mucho antes que la fuerza lateral. (b) El avance
neumático parte de aproximadamente un tercio de la semilongitud de huella y se
anula en el pico de $\Fy$. El modelo de cepillo satura en lugar de caer: un
neumático real pierde fuerza pasado el pico, como en la
\cref{fig:neumatico:curvas}.}
\label{fig:neumatico:mz}
\end{figure}

% ---------------------------------------------------------------------
\section{Elipse de fricción}
\label{sec:neumatico:03}
\index{elipse de fricción}

Hasta aquí se han tratado por separado la fuerza lateral y la longitudinal. En
el coche real nunca lo están: se frena entrando en curva y se acelera saliendo.
Y el neumático tiene \textbf{un solo presupuesto de rozamiento} que repartir
entre las dos.

La aproximación habitual es la elipse de fricción:

\begin{equation}
  \left(\frac{\Fx}{\Fx^{\max}}\right)^{2} +
  \left(\frac{\Fy}{\Fy^{\max}}\right)^{2} \le 1
  \label{eq:elipse}
\end{equation}

La consecuencia práctica está en el panel (b) de la
\cref{fig:neumatico:elipse}, y no es intuitiva: gastar la mitad del agarre
longitudinal solo cuesta un \SI{13}{\percent} del lateral, pero gastar el
\SI{90}{\percent} deja apenas el \SI{44}{\percent}. Es decir, \textbf{se puede
frenar bastante fuerte mientras se gira, pero los últimos gramos de frenada
salen carísimos}. Ahí está la técnica del \emph{trail braking} y ahí está
también la diferencia entre un piloto rápido y uno que no lo es: el lento frena
recto, suelta el freno, gira y luego acelera ---recorriendo los ejes de la
elipse--- mientras que el rápido recorre el borde.

\begin{figure}[htbp]
\centering
\begin{tikzpicture}
\begin{axis}[cursopanel,
  xlabel={$\Fy/\Fz$}, ylabel={$\Fx/\Fz$},
  xmin=-2, xmax=2, ymin=-2, ymax=2,
  xtick={-1.5,-1,...,1.5}, ytick={-1.5,-1,...,1.5},
  title={(a) Presupuesto de agarre}]
% elipse límite
\addplot[curvaA,domain=0:360,samples=200,variable=\t]
  ({1.62*sin(\t)},{1.70*cos(\t)});
% elipses interiores
\addplot[fusalGris,densely dotted,domain=0:360,samples=120,variable=\t]
  ({0.81*sin(\t)},{0.85*cos(\t)});
\addplot[fusalGris,densely dotted,domain=0:360,samples=120,variable=\t]
  ({1.215*sin(\t)},{1.275*cos(\t)});
% punto de trabajo: 70 % de frenada
\addplot[only marks,mark=*,mark size=1.6pt,fusalSec]
  coordinates {(1.157,-1.19)};
\node[cursonota,anchor=north east,text=fusalSec] at (axis cs:1.05,-1.30)
     {\SI{70}{\percent} freno\\\SI{71}{\percent} lateral};
\node[cursonota,anchor=south] at (axis cs:-1.0,1.5) {tracción};
\node[cursonota,anchor=north] at (axis cs:-1.0,-1.5) {frenada};
\end{axis}
\end{tikzpicture}%
\hfill
\begin{tikzpicture}
\begin{axis}[cursopanel,
  xlabel={$\Fx$ usado (\si{\percent} del máximo)},
  ylabel={$\Fy$ disponible (\si{\percent})},
  xmin=0, xmax=100, ymin=0, ymax=105,
  xtick={0,20,...,100}, ytick={0,20,...,100},
  title={(b) Lo que queda para girar}]
\addplot[curvaB,domain=0:100,samples=120]{100*sqrt(1-(x/100)^2)};
\draw[cursomarca] (axis cs:50,0) -- (axis cs:50,86.6) -- (axis cs:0,86.6);
\draw[cursomarca] (axis cs:90,0) -- (axis cs:90,43.6) -- (axis cs:0,43.6);
\node[cursonota,anchor=south west] at (axis cs:52,88) {87};
\node[cursonota,anchor=south west] at (axis cs:60,45) {44};
\end{axis}
\end{tikzpicture}
\caption[Elipse de fricción]{Elipse de fricción de un neumático a carga
constante. (a) La frontera es el límite de lo que el neumático puede dar; el
interior es lo que el piloto está usando. (b) La misma información en la forma
en que se usa para diseñar: fracción de fuerza lateral que queda disponible
según cuánta longitudinal se esté gastando. Curvas sintéticas.}
\label{fig:neumatico:elipse}
\end{figure}

\begin{ojo}
La elipse es una \textbf{aproximación cómoda, no una ley}. En un neumático real
$\Fx^{\max}$ y $\Fy^{\max}$ no coinciden, la frontera no es exactamente
elíptica, y la combinación depende además de la caída y de la carga. Para
dimensionar frenos o repartir par entre ejes es más que suficiente; para un
simulador de vuelta que se quiera creer, hay que usar el modelo combinado del
capítulo \ref{ch:modelado_neumatico}.

El diagrama equivalente a nivel de \emph{coche} ---el diagrama g-g--- no es
simétrico: la frenada llega mucho más lejos que la aceleración, que está
limitada por la potencia y por la tracción del eje motriz. Ese diagrama es el
que sale del capítulo \ref{ch:simulacion_vuelta}, y no debe confundirse con
este.
\end{ojo}

% ---------------------------------------------------------------------
\section{Sensibilidad a la carga vertical}
\label{sec:neumatico:04}
\index{sensibilidad a la carga}\index{transferencia de carga}

Este es, con diferencia, el apartado más rentable del capítulo: explica por qué
la transferencia de carga es mala, por qué las barras estabilizadoras cambian el
balance del coche y por qué bajar el centro de gravedad da más de lo que parece.

\subsection{El agarre no es proporcional a la carga}

Si se dibuja el pico de fuerza lateral frente a la carga vertical, no sale una
recta: sale una curva que se dobla (\cref{fig:neumatico:carga}a). Dicho de otra
forma, el coeficiente de agarre \emph{baja} al aumentar la carga
(\cref{fig:neumatico:carga}b). En el rango de un coche de Formula Student la
caída se aproxima bien por una recta:

\begin{equation}
  \muy(\Fz) \approx \mu_0 - k\,\Fz
  \label{eq:sensibilidad_carga}
\end{equation}

con valores del orden de $\mu_0 \approx \num{1.9}$ y
$k \approx \SI{2.8e-4}{\per\newton}$ para el juego de curvas de este capítulo.

\subsection{Lo que cuesta la transferencia de carga}

Considérese un eje con carga total $2\prom{\Fz}$ repartida entre las dos ruedas,
y una transferencia de carga $\Delta\Fz$ que descarga la interior y carga la
exterior. Sustituyendo \eqref{eq:sensibilidad_carga} en la suma de las dos
ruedas:

\begin{align}
  \Fy^{\text{eje}}
  &= \muy(\prom{\Fz}-\Delta\Fz)(\prom{\Fz}-\Delta\Fz)
   + \muy(\prom{\Fz}+\Delta\Fz)(\prom{\Fz}+\Delta\Fz) \notag \\
  &= \underbrace{2\,\prom{\Fz}\,\bigl(\mu_0 - k\,\prom{\Fz}\bigr)}
     _{\text{sin transferencia}} \;-\; 2\,k\,\Delta\Fz^{2}
  \label{eq:perdida_transferencia}
\end{align}

\begin{clave}
La pérdida de agarre de un eje crece con el \textbf{cuadrado} de la
transferencia de carga, y no depende de cómo se reparta entre muelles, barras o
geometría: solo de cuánta hay. De aquí salen las tres consecuencias que gobiernan
el diseño del coche:

\begin{enumerate}
  \item \textbf{Toda transferencia de carga resta agarre.} Solo se reduce
        bajando el centro de gravedad, ensanchando la vía o quitando masa
        ---nunca tocando muelles o barras---.
  \item \textbf{El reparto entre ejes no cambia el total, pero sí el balance.}
        El eje que se lleva más transferencia pierde más agarre, y por eso el
        $\LLTD$ es la herramienta principal de ajuste de subviraje
        (capítulo \ref{ch:regimen_estacionario}).
  \item \textbf{Quitar masa vale doble.} Menos masa es menos carga vertical
        ---y por tanto más $\muy$--- y además menos transferencia.
\end{enumerate}
\end{clave}

Con los números de la figura y un eje cargado a \SI{1000}{\newton} por rueda, una
transferencia de \SI{500}{\newton} cuesta un \SI{4}{\percent} del agarre del eje;
una de \SI{1000}{\newton} ---justo la que levanta la rueda interior--- cuesta un
\SI{17}{\percent}. Merece la pena tener ese número a mano cuando alguien proponga
subir el radiador \SI{50}{\milli\meter} para que quepa mejor.

\begin{figure}[htbp]
\centering
\begin{tikzpicture}
\begin{axis}[cursopanel,
  xlabel={$\Fz$ (\si{\newton})}, ylabel={$\Fy^{\max}$ (\si{\newton})},
  xmin=0, xmax=2400, ymin=0, ymax=3200,
  xtick={0,500,...,2000}, ytick={0,500,...,3000},
  title={(a) Pico de fuerza lateral}]
\addplot[curvaA,domain=0:2400,samples=80]{x*(1.90-0.00028*x)};
\addplot[fusalGris,densely dotted,domain=0:2400,samples=2]{1.76*x};
\node[cursonota,anchor=west,rotate=42] at (axis cs:1250,2500)
     {proporcional};
\node[cursonota,anchor=north west] at (axis cs:1450,2100) {real};
\end{axis}
\end{tikzpicture}%
\hfill
\begin{tikzpicture}
\begin{axis}[cursopanel,
  xlabel={$\Delta\Fz/\prom{\Fz}$}, ylabel={$\Fy^{\text{eje}}$ (\si{\percent})},
  xmin=0, xmax=1, ymin=80, ymax=101,
  xtick={0,0.2,...,1}, ytick={80,85,...,100},
  title={(b) Coste de la transferencia}]
\addplot[curvaB,domain=0:1,samples=80]{100*(1-0.1728*x^2)};
\draw[cursomarca] (axis cs:0.5,80) -- (axis cs:0.5,95.7) -- (axis cs:0,95.7);
\draw[cursomarca] (axis cs:1,80) -- (axis cs:1,82.7) -- (axis cs:0,82.7);
\node[cursonota,anchor=north east] at (axis cs:0.97,88.5)
     {rueda interior\\en el aire};
\end{axis}
\end{tikzpicture}
\caption[Sensibilidad a la carga vertical]{Sensibilidad a la carga vertical.
(a) El pico de fuerza lateral crece menos que proporcionalmente con la carga: la
recta punteada es lo que daría un rozamiento de Coulomb ideal. (b) Consecuencia
directa: la fuerza lateral que puede dar un eje cae con el cuadrado de la
transferencia de carga, referida a la carga media por rueda. Curvas sintéticas.}
\label{fig:neumatico:carga}
\end{figure}

\begin{ojo}
La sensibilidad a la carga es también la razón por la que \textbf{la carga
aerodinámica rinde cada vez menos}: cada newton de carga vertical adicional
llega con un $\muy$ algo menor que el anterior. En el capítulo
\ref{ch:merece_la_pena_aero} eso se convierte en números; aquí basta con no
suponer que el agarre y la carga vertical van de la mano.
\end{ojo}

% ---------------------------------------------------------------------
\section{Caída, presión y temperatura}
\label{sec:neumatico:05}
\index{caída}\index{presión de inflado}\index{temperatura del neumático}

Los tres parámetros de este apartado tienen algo en común: son baratos, se
cambian en el pit lane en cinco minutos y mueven el agarre más que la mitad de
las decisiones de diseño que ocupan la temporada entera.

\subsection{Caída}

Una caída negativa moderada aumenta la fuerza lateral por dos motivos: aparece un
empuje de caída (\emph{camber thrust}) que se suma al de deriva, y compensa el
vuelco de la huella cuando el coche balancea y el neumático se deforma. Pero hay
un óptimo (\cref{fig:neumatico:cpt}a): pasado ese punto, la huella apoya de
canto, el hombro interior se sobrecarga y se recalienta, y la fuerza cae.

\begin{ojo}
La caída que importa es la \textbf{dinámica en la huella}, con el coche
balanceado, comprimido y con el neumático deformado, no la estática que se mide
en el taller con el coche parado. Un coche con \SI{-2}{\degree} estáticos puede
estar trabajando a \SI{+0.5}{\degree} en el exterior de la curva si la
cinemática no acompaña. Ese es justamente el trabajo del capítulo
\ref{ch:cinematica_suspension}, y la razón de que la curva de ganancia de caída
sea un objetivo de diseño y no una consecuencia.
\end{ojo}

\subsection{Presión de inflado}

La presión cambia a la vez la forma de la huella, la rigidez de deriva, la
rigidez vertical del neumático y el coeficiente de agarre, y todas tienen su
óptimo en sitios distintos:

\begin{itemize}
  \item \textbf{Poca presión}: huella más larga y más flexible, hombros
        sobrecargados, más deformación y más temperatura, respuesta perezosa.
  \item \textbf{Mucha presión}: huella corta y abombada en el centro, menos
        superficie efectiva, respuesta más rápida y menos agarre de pico.
\end{itemize}

Además, la presión sube al calentarse el neumático ---típicamente entre
\num{0.2} y \SI{0.4}{\bar} entre frío y caliente---, así que lo que se ajusta en
el taller no es lo que hay en pista. Lo correcto es fijar la \textbf{presión en
caliente} objetivo y calcular la de frío hacia atrás con los datos de los
primeros entrenamientos.

\begin{clave}
La presión es la variable de puesta a punto \textbf{con mejor relación entre lo
que cuesta y lo que da}: es gratis, se cambia en un minuto y mueve el agarre
varios puntos porcentuales. Si el equipo solo puede permitirse instrumentar y
documentar una variable de setup, que sea esta. Y se mide con un manómetro
bueno: los de gasolinera tienen un error del orden de la magnitud que se intenta
ajustar.
\end{clave}

\begin{regla}[neumáticos]
El reglamento limita lo que se puede hacer con los neumáticos, y algunas de esas
limitaciones condicionan directamente este capítulo: se declaran los juegos de
seco y de lluvia y hay que competir con los declarados; los neumáticos de lluvia
tienen requisitos propios de dibujo o de compuesto; y \textbf{el
acondicionamiento térmico y químico ---mantas térmicas, calentadores,
tratamientos de la goma--- suele estar prohibido}. Eso último no es un detalle
administrativo: significa que el coche tiene que funcionar con el neumático frío
en la primera vuelta, y que la ventana de temperatura del apartado siguiente hay
que alcanzarla conduciendo.

Comprueba los artículos correspondientes del reglamento vigente antes de decidir
neumático, y anota la versión que has consultado: cambian entre temporadas.
\end{regla}

\subsection{Temperatura}

El agarre tiene una ventana (\cref{fig:neumatico:cpt}c). Fuera de ella el
neumático da bastante menos, y en Formula Student casi siempre se está por
debajo: las mangas son cortas, no hay mantas térmicas y las primeras vueltas del
autocross se corren con el neumático frío.

El diagnóstico se hace con un pirómetro de aguja, midiendo tres puntos por
neumático ---interior, centro y exterior--- nada más entrar en boxes:

\begin{center}
\begin{tabularx}{\textwidth}{@{}lX@{}}
\toprule
\textbf{Lo que mide el pirómetro} & \textbf{Qué significa} \\
\midrule
Interior mucho más caliente que exterior & Demasiada caída negativa \\
Exterior mucho más caliente que interior & Falta caída negativa, o el neumático vuelca \\
Centro más caliente que los dos hombros & Presión demasiado alta \\
Hombros más calientes que el centro & Presión demasiado baja \\
Todo el neumático frío & No se le está pidiendo trabajo, o el reparto está mal \\
Un eje mucho más caliente que el otro & Reparto de transferencia desequilibrado \\
\bottomrule
\end{tabularx}
\end{center}

\begin{ojo}
El pirómetro de infrarrojos mide la temperatura \emph{superficial}, que se enfría
en segundos: entre parar el coche y medir se pierde buena parte de la
información. El de aguja llega al interior de la goma y es el que sirve para
comparar entre tandas. Y hay que medir siempre en el mismo orden y a la misma
distancia del box; si no, los números no son comparables entre sesiones.
\end{ojo}

\begin{figure}[htbp]
\centering
\begin{tikzpicture}
\begin{axis}[cursotercio,
  xlabel={$\slipa$ (\si{\degree})}, ylabel={$\Fy$ (\si{\newton})},
  xmin=0, xmax=10, ymin=0, ymax=2000,
  xtick={0,2,...,10}, ytick={0,500,...,2000},
  title={(a) Caída}, legend pos=south east]
\addplot[curvaA,domain=0:10,samples=100]{mf(x,0.3226,1.5,1620,0.3)};
  \addlegendentry{$\camber=0$}
\addplot[curvaB,domain=0:10,samples=100]{mf(x+0.4,0.3226,1.5,1700,0.3)};
  \addlegendentry{$\camber=\SI{-2}{\degree}$}
\addplot[curvaC,domain=0:10,samples=100]{mf(x+0.7,0.3226,1.5,1560,0.3)};
  \addlegendentry{$\camber=\SI{-4}{\degree}$}
\end{axis}
\end{tikzpicture}%
\hfill
\begin{tikzpicture}
\begin{axis}[cursotercio,
  xlabel={$\pneu$ (\si{\bar})}, ylabel={$\muy$ de pico},
  xmin=0.5, xmax=1.2, ymin=1.45, ymax=1.68,
  xtick={0.5,0.7,...,1.1}, ytick={1.45,1.50,...,1.65},
  title={(b) Presión}]
\addplot[curvaB,domain=0.5:1.2,samples=80]{1.62-1.0*(x-0.83)^2};
\draw[cursomarca] (axis cs:0.83,1.45) -- (axis cs:0.83,1.62);
\node[cursonota,anchor=south] at (axis cs:0.95,1.63) {óptimo};
\end{axis}
\end{tikzpicture}%
\hfill
\begin{tikzpicture}
\begin{axis}[cursotercio,
  xlabel={$\Tneu$ (\si{\celsius})}, ylabel={$\muy/\muy^{\max}$},
  xmin=30, xmax=130, ymin=0.5, ymax=1.05,
  xtick={30,50,...,130}, ytick={0.5,0.6,...,1.0},
  title={(c) Temperatura}]
\addplot[fusalPrim!12,draw=none,domain=65:95,samples=2,
         fill=fusalPrim!12] {1} \closedcycle;
\addplot[curvaC,domain=30:130,samples=100]{1-((x-78)/85)^2};
\node[cursonota,anchor=south] at (axis cs:80,0.55) {ventana\\de trabajo};
\end{axis}
\end{tikzpicture}
\caption[Efecto de caída, presión y temperatura]{Los tres parámetros que se
tocan en el pit lane. (a) Una caída negativa moderada sube la fuerza lateral;
pasado el óptimo la pierde. (b) El agarre de pico tiene un óptimo de presión.
(c) La goma solo da su máximo dentro de una ventana de temperatura. Los tres
paneles son \textbf{cualitativos}: el óptimo concreto es distinto para cada
neumático y cada compuesto, y solo se encuentra midiendo.}
\label{fig:neumatico:cpt}
\end{figure}

% ---------------------------------------------------------------------
%  Cierre del capítulo
% ---------------------------------------------------------------------

\begin{caso}[Pirelli 185/40 R13]
Montamos un Pirelli \texttt{185/40 R13} de compuesto duro. La decisión no fue de
prestaciones sino de \textbf{flujo de trabajo}: renunciamos a un
\SIrange{5}{10}{\percent} de agarre lateral frente a un slick típico de
Formula Student a cambio de que dure muchísimo más, porque lo que este equipo
necesita ahora mismo no son décimas, son kilómetros de rodaje para instrumentar
el coche y aprender a conducirlo. Quemar neumáticos blandos en entrenamientos no
tiene sentido cuando todavía no se sabe qué se está midiendo.

Consecuencias que hay que asumir y saber defender:

\begin{itemize}
  \item \textbf{No tenemos datos de este neumático} ---no está en el TTC---, así
        que en el modelo tratamos el agarre de forma conservadora. Cómo se hace
        eso y cuánto error se arrastra es el capítulo
        \ref{ch:modelado_neumatico}.
  \item La caída estática ---\SI{-1.5}{\degree} delante y \SI{-1.0}{\degree}
        detrás--- se fijó con la curva de ganancia de la herramienta de
        cinemática y experiencia de otros equipos, \emph{no} con datos de agarre.
        Es una base razonable, no un óptimo.
  \item El soporte de mangueta es intercambiable justamente para poder probar
        otros valores en cuanto podamos medir agarre de verdad.
\end{itemize}

\pendiente{Completar con: presión en frío y en caliente de trabajo, temperaturas
medidas con el pirómetro por tanda, y desgaste observado a lo largo de la
temporada. Sin estos tres datos este recuadro sigue siendo una intención, no un
caso.}
\end{caso}

\begin{ojo}
\begin{itemize}
  \item Usar directamente el $\muy$ de los datos de banco. Sale un coche
        diseñado para un agarre que en pista no existe.
  \item Hablar de «agarre» sin decir a qué carga vertical: sin $\Fz$, un
        coeficiente de agarre no significa nada.
  \item Ajustar la caída estática con el coche parado y no comprobar nunca la
        dinámica.
  \item Inflar en frío al valor que se quiere en caliente.
  \item Medir temperaturas con el infrarrojo tres minutos después de parar y
        sacar conclusiones de esos números.
  \item Suponer que un neumático más ancho agarra proporcionalmente más: gana
        por área de contacto y por menor sensibilidad a la carga, pero también
        pesa más, tiene más inercia y más resistencia.
  \item Cambiar dos cosas a la vez entre tandas y no poder atribuir la mejora a
        ninguna.
\end{itemize}
\end{ojo}

\begin{entregable}
\begin{enumerate}
  \item La \textbf{ficha del neumático} en una página: medida, compuesto, llanta
        recomendada, radio de rodadura, presiones de trabajo en frío y en
        caliente, caída objetivo y ventana de temperatura. Es la primera hoja de
        la carpeta de dinámica vehicular.
  \item La \textbf{hoja de registro de tanda}: presiones en frío y en caliente y
        temperaturas en tres puntos por neumático, con hueco para lo que se
        cambió y qué dijo el piloto. Impresa, con pinza, en el box.
  \item Una \textbf{hoja de cálculo de sensibilidad a la carga} que, dados
        $\mu_0$ y $k$, dé la pérdida de agarre de un eje frente a la
        transferencia de carga según \eqref{eq:perdida_transferencia}.
\end{enumerate}
\end{entregable}

\begin{juez}
\begin{enumerate}
  \item ¿Por qué habéis elegido este neumático y qué habéis sacrificado?
  \item ¿Cuál es el coeficiente de agarre que usáis en vuestros cálculos y de
        dónde sale?
  \item ¿Cuánta transferencia de carga tenéis en el eje delantero a
        \SI{1.4}{\gacc} y cuánto agarre os cuesta?
  \item ¿Cuál es la caída dinámica en la rueda exterior a máximo balanceo?
  \item ¿A qué presión rodáis y cómo habéis llegado a ese número?
  \item ¿Qué temperatura alcanzan vuestros neumáticos y cómo lo sabéis?
\end{enumerate}
\end{juez}
